
\chapter{Споры и~возвратные платежи}\label{ch:disputes}

\highlight{Когда держатель карты звонит в~банк и~говорит: \enquote{я~этого
не~покупал}, вчерашняя завершённая транзакция превращается в~спор}. Каждая
стадия меняет денежную позицию сторон, бремя доказывания и~срок ответа.

\section{Пять стадий диспута}\label{sec:disputes-lifecycle}

\highlight{Получив жалобу, банк запускает регламентированный пятистадийный
  цикл; каждая стадия задаёт срок, набор документов и~финансовые
последствия}~\cite{visa-core-rules,mc-chargeback-guide}.

Первая стадия~--- предварительный запрос на~сведения, в~отрасли~---
\textit{inquiry} или \textit{retrieval request}: эмитент или его процессор
проверяет, есть ли основания открывать формальный спор. Стадия не~обязательна,
но часть недоразумений снимается ещё до~возвратного платежа.
Если спор не~снят, начинается формальный возвратный платёж без согласия
продавца, \textit{chargeback}. Эмитент обращается к~сети, сеть дебетует
эквайера, а~эквайер переносит риск на~ТСП. К~спору прикладывается код основания
спора, \textit{reason code}: он задаёт предмет доказывания и~показывает, какая
сторона несёт основное бремя доказывания~\cite{visa-core-rules}.

Следующая стадия~--- ответ ТСП на~возвратный платёж, в~отраслевой речи
повторное представление, \textit{representment}. У~Visa официальная
терминология устроена иначе: \enquote{диспут / ответ на~диспут / предарбитраж /
арбитраж}; Mastercard сохраняет язык циклов возвратного платежа и~подачи дела.

Если ответ не~убедил эмитента, наступает предарбитраж,
\textit{pre-arbitration}. Эмитент либо принимает ответ, либо переводит спор
на~более узкую и~дорогую стадию. Если предарбитраж не~закрывает спор, сеть
выносит дело в~арбитраж, \textit{arbitration}, где решение становится
обязательным, а~проигравшая сторона платит по~спорной операции
и~по~процессуальным сборам~\cite{visa-core-rules,mc-chargeback-guide}.

Логи, чеки, сведения о~доставке и~согласии клиента хранят дольше
самого заказа: код основания спора, под который их потребуют, сеть выберет
позже, чем продавец закроет заказ. \FloatBarrier

\section{Процедуры Visa и~Mastercard}\label{sec:disputes-vcr-mc}

В~2018 году Visa перевела процесс споров на~\emph{Visa Claims Resolution}
(VCR). Стороны подают доказательства раньше, а~количество обменов между
эмитентом и~эквайером сокращено~\cite{visa-vcr-merchants,visa-core-rules}.

Когда клиент не~узнаёт списание и~звонит в~банк, эмитент вместо формального
возвратного платежа сначала показывает ему дополнительные сведения о~покупке.
Понятное описание товара, адрес доставки или имя получателя снимают часть
споров до~формального этапа. На~этом построен Order Insight~--- сервис Visa
(через сеть Verifi), передающий эмитенту расширенные данные о~покупке
до~формального возвратного платежа, чтобы закрыть спор до~перехода в~дорогую
процессуальную стадию~\cite{visa-order-insight}.

У~Mastercard задача та~же, но язык правил другой. Официальное руководство говорит
о~циклах возвратного платежа, втором предъявлении (\textit{second presentment})
и~подаче дела в~арбитраж~\cite{mc-chargeback-guide}. За~прояснение покупки
до~формального спора у~Mastercard отвечает семейство сервисов Ethoca для раннего
обмена данными между эмитентом и~ТСП, включая Consumer Clarity, которая передаёт
эмитенту обогащённое описание покупки в~момент звонка
держателя~\cite{mc-dispute-management,mc-ethoca-consumer-clarity}.

Самое заметное процессуальное расхождение~--- сроки. Visa последовательно сокращала
срок ответа на~возвратный платёж: с~45~дней до~30, затем до~20, а~с~21~июля
2025~года~--- до~9~дней для US/Canada и~18~дней для остальных
регионов~\cite{visa-core-rules}. Mastercard сохраняет 45~дней на~ответ эквайера
в~первом цикле возвратного платежа~\cite{mc-chargeback-guide}. У~ТСП из-за
этого складываются разные внутренние сроки на~ответ (Service Level Agreement,
SLA) под каждую сеть: спор по~Visa требует ответа быстрее, чем аналогичный
спор по~Mastercard. Точный срок сверяйте с~текущей редакцией правил.

Исторически 45~дней соответствовали повседневной работе 1970-х--1980-х годов.
Международная почта с~подтверждением доставки шла несколько недель, банковская
обработка добавляла ещё несколько дней. Более длинные сроки для поздних стадий,
предарбитража и~арбитража, давали сторонам время собрать доказательства через
несколько часовых поясов. Сокращение Visa до~9/18 дней отражает переход
на~цифровую основу: логи, отметки курьерской службы, данные 3-D Secure (3DS,
см.~главу о~3-D Secure) доступны немедленно.

Для раннего снятия спора у~Mastercard есть отдельный сервис Ethoca Alerts:
эмитент уведомляет ТСП о~поданной жалобе, и~продавец успевает провести возврат
до~формального возвратного платежа~\cite{mc-dispute-management}.

\begin{figure}[htbp]
  \centering
  \includefiguresvg[width=\linewidth]{assets/figures/ch31-disputes/vcr-vs-mc-flow}
  \caption{Стадии разрешения спора у~Visa и~Mastercard.}\label{fig:vcr-vs-mc-flow}
\end{figure}

Раннее снятие спора применимо только там, где спор возник из-за неузнавания
покупки или недостатка сведений. При неполученном товаре, дефектной услуге или
нарушенном условии договора спор всё равно придётся вести по~полной схеме.

К~Order Insight у~Visa и~к~Ethoca Consumer Clarity у~Mastercard ТСП
подключается через своего эквайера или напрямую через одобренного партнёра
сети~--- Verifi у~Visa, Ethoca у~Mastercard. Оба продукта не~меняют
авторизационный поток: данные отдаются эмитенту по~отдельному API уже после
авторизации, в~момент, когда держатель карты интересуется покупкой или подаёт
жалобу. У~Visa запрос обслуживается через Visa Resolve Online (VROL),
к~которому подключена сеть Verifi; у~Mastercard~--- через Mastercom
и~партнёрскую сеть Ethoca. Для ТСП интеграция сводится к~CRM-коннектору,
который по~входящему запросу возвращает расширенное описание
покупки~--- название позиции, адрес доставки, имя получателя, идентификатор
заказа на~стороне ТСП. Эмитент показывает эти данные держателю карты при звонке
в~колл-центр или при просмотре деталей операции в~мобильном или интернет-банке.
ТСП может вести собственный лог звонков и~возвратов, но без подключения
к~Verifi или Ethoca эти данные до~эмитента не~доходят. \FloatBarrier

\section{Коды основания спора: четыре категории Visa}\label{sec:disputes-reason-codes}

\highlight{Visa делит коды основания спора на~четыре большие категории: фрод
  (10.x), ошибки авторизации (11.x), ошибки обработки (12.x) и~потребительские
споры (13.x)}~\cite{visa-core-rules}. Категория сразу определяет, кто несёт
бремя доказательства и~какие доказательства сеть готова рассматривать.

Категория 10.x~--- фрод (\textit{fraud}). Для интернет-торговли определяющим
остаётся код 10.4, \textit{Other Fraud --- Card-Absent Environment}: держатель
карты утверждает, что не~участвовал в~интернет-покупке. Под код 10.4 у~Visa
и~работает Compelling Evidence~3.0 (CE~3.0;
см.~раздел~\ref{sec:disputes-ce30}).

Категория 11.x~--- нарушение правил авторизации: операция без обязательного
запроса авторизации, продолжение несмотря на~отказ, проведение по~устаревшему
ответу авторизации.

Категория 12.x~--- ошибки обработки на~стороне продавца или процессинга: сумма
клиринга не~совпала с~авторизацией, одна и~та~же операция попала в~клиринг
дважды, транзакция дошла до~сети слишком поздно или с~искажёнными данными.

Категория 13.x описывает потребительские споры: товар не~получен, услуга
не~оказана, подписка не~отменена, возврат обещан, но не~проведён. Для этой
категории решает доказательство исполнения обязательства, а~не технический лог:
доставка, акцепт условий, подтверждение пользования, дата отмены и~момент
начисления следующего биллингового цикла.

У~Mastercard собственная нумерация кодов основания спора, но деление сходное:
фрод, авторизация, ошибка обработки, спор держателя
карты~\cite{mc-chargeback-guide}.

По~коду 10.4 бремя доказательства фрода в~дистанционном канале несёт ТСП.
Поэтому его защита опирается на~3DS, профилирование устройств и~признаки
поведенческого совпадения. По~кодам~13.x в~потребительских спорах предмет
доказывания смещается к~исполнению обязательства: подтверждению доставки
и~раскрытию условий до~транзакции.

У~подписочных кодов доказательная база делится по~существу спора. По~коду 13.2
\enquote{Cancelled Recurring Transaction} Visa с~октября 2024~года принимает
как защиту ТСП доказательство того, что клиент пользовался услугой уже после
даты отзыва разрешения на~списание и~до~даты обработки диспута, либо
документально подтверждённый возврат или зачёт списанной
суммы~\cite{visa-core-rules}. Сценарий \enquote{отписался заранее, а~деньги
списали} проходит по~13.2: держатель отозвал согласие, а~списание прошло
позже. Сценарий \enquote{не~знал о~продлении подписки}~--- триал или
промопериод без~явного уведомления~--- проходит по~коду
13.5 \enquote{Misrepresentation}: здесь Visa с~2020~года требует от~ТСП
подтверждения явного согласия клиента на~будущие списания и~напоминания
не~позднее чем за~7~дней до~окончания пробного или промопериода
(см.~раздел~\ref{sec:subscriptions-soft-decline})~\cite{visa-core-rules,visa-scf}.
У~Mastercard для подписок с~биллингом реже одного раза в~полгода действует
своё правило: за~7--30~дней до~следующего списания клиенту отправляется
напоминание о~предстоящем платеже~\cite{mc-tpr}. Лог коммуникаций ТСП с~датами
отправок и~подтверждениями доставки нужен в~обоих сценариях наравне
с~подтверждением исполнения услуги.

\section{Состав доказательств}\label{sec:disputes-evidence}

Выражение \enquote{убедительные доказательства} (compelling evidence) в~споре
означает конкретный набор фактов под конкретный код основания
спора~\cite{visa-core-rules}. Факты, достаточные для спора о~фроде,
не~сработают в~споре о~неполученном товаре.

В~споре о~дистанционной покупке ТСП показывает привязку транзакции
и~к~устройству, и~к~поведению клиента. Сюда относятся IP-адрес, идентификатор
или отпечаток устройства, учётная запись, адрес доставки, адрес электронной
почты, номер телефона, а~при наличии 3-D Secure результат аутентификации. Один
IP-адрес сам по~себе слаб; убедительной становится совокупность признаков.

Для спора о~неполучении товара или услуги требуется подтверждение исполнения:
отметка курьерской службы о~вручении, подпись получателя, технический лог
активации цифрового товара, запись о~первом входе в~сервис, протокол оказания
услуги. Для спора об~отменённой подписке требуются дата отмены, условия договора
и~доказательство того, что клиент видел эти условия до~списания.

Логи под спор собирайте до~возвратного платежа, а~не во~время него: IP-адрес,
строку \mbox{\texttt{User-Agent}}, сведения об~устройстве, результат проверки адреса
(Address Verification System, AVS) и~кода проверки карты (Card Verification
Value, CVV), результат 3-D Secure, подтверждение доставки и~связку между первой
клиентской транзакцией (Customer-Initiated Transaction, CIT; см.~главу
о~хранимых учётных данных) и~последующими списаниями.

\subsection*{Что не считается доказательством: SMS клиента}\label{sec:disputes-sms-not-evidence}

Держатель часто предъявляет как \enquote{доказательство} SMS-уведомление
от~банка-эмитента~--- \enquote{с~карты списано 1500~рублей в~ТСП X}. Когда чек
продавцу не~достался или терминал остановился на~печати,
держатель карты ссылается на~это сообщение как на~подтверждение транзакции.
В~цикле спора оно аргументом не~является: ни одна из~карточных сетей
не~специфицирует SMS-информирование и~не учитывает его в~правилах разрешения
спора~\cite{visa-core-rules,mc-tpr}.

Само обязательство SMS-информирования вытекает из~части~4 статьи~9
Федерального закона от~27.06.2011~№~161-ФЗ~\cite{fz-161}: эмитент обязан
уведомлять клиента о~каждой операции с~использованием его электронного средства
платежа (ЭСП; см.~главу о~161-ФЗ), способ уведомления согласуется в~договоре,
и~с~момента уведомления начинается отсчёт срока, в~течение которого клиент
вправе оспорить операцию у~эмитента. Положение Банка России
от~17.08.2023~№~821-П~\cite{cbr-821p} регулирует защиту канала:
подтверждение контактных данных клиента и~правила обработки информации
о~переводах. Межбанковский и~клиентский циклы спора по~срокам пересекаются лишь
частично.

SMS отправляет внутренний сервис эмитента по~триггеру изменения доступного
остатка, формируя текст из~полей сообщения ISO~8583 (DE~--- Data Element;
см.~главу о~ISO~8583): DE~41 (Terminal ID), DE~42 (Merchant ID) и~DE~43 (Card
Acceptor Name/Location) авторизационного сообщения. Триггер срабатывает
на~одобренной авторизации, но не~подтверждает, что ответ хоста дошёл
до~терминала, был корректно разобран POS-терминалом (Point-of-Sale)
и~распечатался чек. Терминал мог выполнить автоматический технический реверс
(\textit{reversal}) по~таймауту или по~ошибке разбора криптограммы ответа
авторизации (Authorization Response Cryptogram, ARPC; см.~главу о~EMV) уже
после отправки SMS держателю. Держатель карты видит \enquote{списание},
продавец~--- отказ.

В~доказательную базу спора входит только то, что прошло через клиринг
и~карточную сеть: со~стороны держателя~--- чек терминала с~кодом ответа DE~39
(Response Code), со~стороны продавца~--- клиринговая запись с~тем же кодом
авторизации (\textit{auth code}). Push-уведомления мобильного банка
и~отображение в~истории операций до~клиринга сеть не~рассматривает.

\section{Compelling Evidence 3.0 и~NTID}\label{sec:disputes-ce30}

С~15~апреля 2023~года Visa расширила правила доказательства и~ввела Compelling
Evidence~3.0 (CE~3.0) для защиты от~спора о~фроде в~дистанционном
канале при наличии двух предшествующих неоспоренных транзакций с~тем же
реквизитом и~совпадающими ключевыми признаками клиента. Так Visa отвечает
на~дружественный фрод (\textit{friendly fraud};
см.~раздел~\ref{sec:fraud-threats})~--- ситуацию, когда легитимный держатель
оспаривает собственную покупку как
несанкционированную~\cite{visa-ce30-merchant-readiness,visa-core-rules}. Связку
списаний на~всём интервале обеспечивает \emph{сетевой идентификатор транзакции}
(Network Transaction ID, NTID)~--- присваиваемый сетью ключ, привязывающий
последующие операции к~первой клиентской.

Для защиты по~CE~3.0 требуются две предыдущие транзакции того же ТСП, у~которых
нет ни активного отчёта о~фроде, ни активного спора по~фроду, и~которые
попадают в~интервал 120--365~дней до~даты спора. Между спорной
и~историческими транзакциями совпадают минимум два ключевых признака клиента,
и~одним из~них обязан быть IP-адрес либо идентификатор
устройства~\cite{visa-ce30-merchant-readiness}.

На~уровне протокола CE~3.0 работает через Visa Resolve Online (VROL)~--- единую
онлайн-платформу Visa для обмена данными споров между эмитентом и~эквайером.
При споре с~\textit{reason code} 10.4 \enquote{Other Fraud --- Card-Absent
Environment} система сама проверяет, есть ли в~истории эквайера подходящие
транзакции, и~при подтверждённых совпадениях отклоняет диспут
(\textit{dispute}) на~стадии \textit{pre-dispute} или \textit{dispute} без
ручного ответа ТСП. Условие: заранее передавать device ID и~IP в~данных
авторизации и~хранить историю неоспоренных транзакций минимум
365~дней~\cite{visa-ce30-merchant-readiness}.

Visa не~требует NTID как отдельный критерий CE~3.0. Но в~подписочном сценарии
сохранённый NTID и~аккуратная связка первой клиентской транзакции CIT
с~последующими списаниями дают ТСП исторический след
клиента~\cite{visa-ce30-merchant-readiness,visa-scf}.

Mastercard развивает First-Party Trust Program~--- программу сопоставимого
назначения против того же явления
(у~Mastercard используется термин \emph{first-party chargeback}). ТСП может передавать расширенные данные либо
в~момент авторизации, либо на~стадии спора; для защиты от~возвратного платежа
программа требует комбинацию признаков из~групп device factor (факторы
устройства), delivery factor (факторы доставки) и~additional identity factors
(дополнительные идентифицирующие
признаки)~\cite{mc-first-party-trust,mc-dispute-management}.

У~CE~3.0 узкая область действия: только споры по~фроду в~дистанционном канале.
Против неполученного товара, дефектной услуги или ошибочной отмены подписки она
не~работает.

\section{Мониторинговые программы сетей}\label{sec:disputes-monitoring}

Устойчиво высокий уровень возвратных платежей грозит попаданием в~мониторинговые
программы сетей. Названия и~пороги у~программ меняются, но одно не~меняется: при
слишком высокой доле споров или фрода сеть переходит от~предупреждений к~плану
исправления, штрафам и~к~прекращению
обслуживания~\cite{visa-core-rules,mc-chargeback-guide}.

В~публичных материалах Visa свела мониторинг споров и~фрода на~уровне эквайера
в~Visa Acquirer Monitoring Program (VAMP), а~Mastercard
по-прежнему описывает набор отдельных программ контроля, включая Business Risk
Assessment and Mitigation (BRAM), Excessive Chargeback Program (MC~ECP),
Excessive Fraud Merchant (EFM) и~Questionable Merchant Audit Program
(QMAP)~\cite{visa-vamp-intro,mc-compliance-programs}.

VAMP мониторит оба уровня~--- и~портфель эквайера в~целом, и~отдельные ТСП,
с~разными порогами для каждого. Программа
объединила прежние Visa Dispute Monitoring Program (VDMP) и~Visa Fraud
Monitoring Program (VFMP) в~единый мониторинг. В~публичном документе Visa,
действовавшем с~1~июня 2025~года, для портфеля эквайера сохранены уровни Above
Standard и~Excessive, тогда как для ТСП уровень Above Standard был отменён:
остался единый глобальный порог Excessive 220~базисных пунктов при не~менее
1\,500 событий фрода и~диспутов в~месяц; применение мер по~нему началось
с~1~октября 2025~года; Visa отдельно объявила снижение этого порога
до~150~базисных пунктов с~1~апреля 2026~года~\cite{visa-vamp-fact-sheet-2025}.
Точный порог сверяйте с~датой редакции программы. При превышении порога
участник попадает в~мониторинговую программу и~обязан внедрять меры по~снижению
риска.

У~Mastercard программа мониторинга возвратных платежей ТСП оформлена в~MC~ECP.
Публичные Security Rules and Procedures выносят BRAM и~MC~ECP в~разные
программы, а~для ТСП с~превышением порогов по~возвратным платежам используют
категории Excessive Chargeback Merchant (ECM) и~High Excessive Chargeback
Merchant (HECM)~\cite{mc-compliance-programs,mc-spme-2025}. За~конкретными
числовыми порогами Mastercard отсылает участников к~руководству Data Integrity
Monitoring Program. В~публичных материалах MC~ECP описывается как отдельная
программа с~эскалацией у~эквайера, без привязки к~фиксированному
порогу BRAM.
Широко используемые ориентиры: для ECM количество возвратных платежей $\geq$~100
в~месяц при доле возвратных платежей (\textit{chargeback ratio})
$\geq$~1,5\,\%; для HECM~--- $\geq$~300 при доле $\geq$~3,0\,\%. Доля
возвратных платежей рассчитывается как количество возвратных платежей
за~текущий месяц, делённое на~количество \textit{e-commerce}-транзакций
за~предшествующий месяц~\cite{mc-compliance-programs}.

В~плане исправления сеть требует конкретных мер: улучшение дескриптора
продавца (\textit{merchant descriptor}~--- текстовая строка с~названием
и~местоположением ТСП, которая попадает в~выписку держателя карты), подключение
к~Order Insight или Ethoca, ужесточение проверки 3DS, пересмотр политики
возвратов, ограничение проблемных категорий товаров. Без снижения показателей
за~90--120 дней эскалация продолжается.

Пороги \enquote{CB ratio $>$ 0{,}9\,\%} и~штрафы \$25--100 за~возвратный
платёж относятся к~эпохе VDMP, а~не к~действующей редакции VAMP.
Знаменатели у~методик расчёта разные: VAMP считает в~базисных пунктах по~формуле
$(\text{TC}40+\text{TC}15)/\text{TC}05_{\text{CNP}}$, где TC~--- Transaction
Code в~клиринге Visa BASE~II (TC~40~--- отчёт о~фроде, TC~15~--- возвратный
платёж, TC~05~--- продажа; см.~главу о~клиринге).

Сеть оценивает сочетание абсолютного объёма и~относительной доли. Сезонность,
скользящий период и~сдвиг спора в~соседний календарный месяц легко искажают
показатель; ТСП надёжнее считать свои показатели отдельно по~спорам, фроду
и~возвратам, чем смотреть на~единый агрегированный показатель.

\section{Диспуты в~системе \enquote{Мир}}\label{sec:disputes-mir}

Карты \enquote{Мир} обрабатываются ОПКЦ НСПК (раздел~\ref{sec:mir-opkc}),
и~регламент диспутов по~ним задают Правила платёжной системы \enquote{Мир};
Visa Resolve Online и~Mastercom (единая платформа управления спорами Mastercard
для обмена сообщениями и~доказательствами между эмитентом и~эквайером) здесь
не~применяются~\cite{nspk-mir-rules}.\footnote{Конкретные коды оснований, сроки
  и~тарифы стадий диспута публикуются НСПК в~правилах платёжной системы
  \enquote{Мир} для участников; числовые значения и~состав кодов сверяются
с~действующей редакцией перед публикацией.} По~устройству регламент НСПК
повторяет общую модель карточного диспута: возвратный платёж
\textit{chargeback}, повторное представление \textit{representment}, претензия
\textit{pre-arbitration}, арбитраж. Различия проявляются в~кодах оснований,
сроках и~тарифах за~стадии.

Регламент \enquote{Мир} делит основания диспутов на~несколько групп.
Первая~--- операции, оспариваемые держателем карты: не~получил товар или
услугу, неавторизованная операция, дубль операции, неверная сумма. Вторая~---
операции, оспариваемые эмитентом по~техническим основаниям: неверная информация
о~ТСП, нарушение правил процессинга, технические сбои. Третья~--- операции
с~подозрением на~фрод. НСПК устанавливает собственные сроки оспаривания:
максимальный срок подачи диспута для большинства оснований составляет 180~дней
с~даты операции. Тарифы
за~стадии (плата эквайера за~принятие возвратного платежа, плата при повторном
представлении, плата за~арбитраж) фиксированы в~тарифах
ПС~\cite{nspk-mir-tariffs} и~платятся в~рублях.

Для продавца доказательная база и~сроки ответа на~возвратный платёж по~карте
\enquote{Мир} из~привычного руководства по~спорам Visa/Mastercard заимствовать
напрямую нельзя; состав обязательных доказательств и~сроки ответа сверяйте
с~правилами \enquote{Мир} через эквайера и~его процессор. Для эквайера
мониторинговые программы НСПК (аналоги VAMP/MC~ECP) действуют отдельно. Пороги
доли возвратных платежей и~доли фрода, после которых ТСП
попадает под повышенный контроль или санкции, устанавливает НСПК.

\section{Диспуты в~СБП}\label{sec:disputes-sbp}

Держатель карты, недовольный покупкой, нажимает кнопку спора в~банковском
приложении, и~карточный возвратный платёж начинает свой цикл. Плательщик через
СБП в~типовом потоке такой кнопки не~видит: у~системы нет встроенного механизма
принудительного оспаривания уже исполненной операции.

В~СБП деньги зачисляются получателю мгновенно. При ошибочном переводе банк
плательщика может запустить \enquote{Просьбу СБП}, но автоматически вернуть
деньги нельзя~--- банк обращается к~получателю, потому что для возврата
требуется его согласие, если только иной порядок не~вытекает из~закона, решения
суда или договора с~банком~\cite{cbr-sbp-faq-mistaken-transfer}. Для оплаты
бизнеса доступен обычный возврат через банк продавца, в~том числе частичный,
и~деньги при таком возврате поступают покупателю
моментально~\cite{sbp-faq-business}.

Продавец не~несёт классических штрафов по~спорам и~не~проходит цепочку
от~ответа на~возвратный платёж до~предарбитража. Покупатель, если товар
не~доставлен, защищает свои интересы через возврат, который должен оформить
продавец, переговоры, претензионную работу или обычный гражданско-правовой спор.

В~продукте, который принимает и~карты, и~СБП, эти два потока работают по~разным
правилам: карточный~--- с~циклом возвратного платежа и~SLA сети, СБП~---
через гражданско-правовой возврат и~\enquote{Просьбу}.
